\begin{abstract}
We prove the critical $L^3$ local smoothing estimate for the wave equation
in three spatial dimensions, with every positive Sobolev loss. This resolves
Sogge's Euclidean local smoothing conjecture in dimension three.
\end{abstract}

\section{Introduction}\label{sec:introduction}

Let
\[
 \widehat f(\xi)=\int_{\R^3}e^{-ix\cdot\xi}f(x)\,dx,
 \qquad
 Uf(x,t)=\frac1{(2\pi)^3}
       \int_{\R^3}e^{i(x\cdot\xi+t|\xi|)}\widehat f(\xi)\,d\xi.
\]
For $\alpha\in\R$, write $J^\alpha=(1-\Delta)^{\alpha/2}$, with
Fourier multiplier $(1+|\xi|^2)^{\alpha/2}$.
For $p>2$, the sharp fixed-time $L^p$ estimate has Sobolev loss
$2(1/2-1/p)$ in three dimensions \cite{Peral1980,Miyachi1980}.
Time integration can recover part of that loss. Sogge's Euclidean
local smoothing conjecture \cite{Sogge1991} predicts, for $2<p<\infty$,
the strict range $\alpha>\max\{0,1-3/p\}$; its critical exponent is
$p=3$. Our main result treats that exponent.

\begin{theorem}[Critical local smoothing]\label{thm:main}
For every $\eps>0$ there exists $C_\eps<\infty$ such that
\begin{equation}\label{eq:intro-main}
 \norm{Uf}_{L^3(\R^3\times[1,2])}
 \le C_\eps\norm{J^\eps f}_{L^3(\R^3)}
 \qquad (f\in\mathcal S(\R^3)).
\end{equation}
\end{theorem}

The critical estimate supplies the interpolation point for the full
Euclidean range: the $L^2$ energy estimate gives the exponents below
$p=3$, and an annular $L^\infty$ kernel bound gives those above it.
Frequency summation then yields Corollary~\ref{cor:full-range}:
\begin{equation}\label{eq:intro-full}
 \norm{Uf}_{L^p(\R^3\times[1,2])}
 \le C_{p,\alpha}\norm{J^\alpha f}_{L^p(\R^3)},
 \qquad
 2<p<\infty,\quad \alpha>\max\{0,1-3/p\}.
\end{equation}
Together, these estimates resolve Sogge's Euclidean local smoothing
conjecture in three spatial dimensions. At $p=3$, the estimate holds with
every positive Sobolev loss; the lossless estimate $\eps=0$ is not addressed.

There is also a pointwise consequence for the wave evolution. Define the
finite-time maximal half-wave by
$U_*f(x)=\sup_{0<t<1}|Uf(x,t)|$. For $3\le p<\infty$ and
$s>1-2/p$, Corollary~\ref{cor:wave-maximal} gives
$\norm{U_*f}_p\le C_{p,s}\norm{J^s f}_p$ and almost-everywhere
convergence $Uf(x,t)\to f(x)$ as $t\downarrow0$ for data with
$J^s f\in L^p$. The established implication from local smoothing uses
Sobolev embedding in time. Its extra cost explains the stronger
regularity hypothesis: at $p=3$, spacetime smoothing allows every
positive loss, whereas this maximal conclusion requires $s>1/3$.

The local-smoothing estimate also has a consequence for Fourier summation.
For $R>0$, $\delta>0$, and $f\in\mathcal S(\R^3)$, define the
Bochner--Riesz means and their maximal operator by
\begin{equation}\label{eq:intro-br-means}
 \begin{split}
 B_R^\delta f(x)
 &=\frac1{(2\pi)^3}\int_{\R^3}e^{ix\cdot\xi}
       \left(1-\frac{|\xi|^2}{R^2}\right)_+^\delta
       \widehat f(\xi)\,d\xi,\\
 B_*^\delta f(x)&=\sup_{R>0}|B_R^\delta f(x)|.
 \end{split}
\end{equation}
Propositions~3.2 and~3.3 of Beltran--Hickman--Sogge
\cite{BeltranHickmanSogge2019Survey} turn local smoothing into
Bochner--Riesz estimates. Section~\ref{sec:bochner-riesz} applies them to
\eqref{eq:intro-full}: Theorem~\ref{thm:br-maximal} bounds $B_*^\delta$
strongly on $L^p$ and gives $B_R^\delta f\to f$ almost everywhere for
$3\le p<\infty$ and $\delta>1-3/p$. Separately,
Corollary~\ref{cor:br-nonmaximal} gives uniform bounds for the individual
means throughout the strict range
$\delta>\max\{3|1/p-1/2|-1/2,0\}$, $1\le p\le\infty$.
The individual bounds control Fourier summation uniformly in its
radius. The stronger maximal bound controls all radii simultaneously
and makes convergence stable under $L^p$ approximation; this is what
permits almost-everywhere recovery of arbitrary $L^p$ data. Here the
strict condition on the summation order absorbs the positive loss in
local smoothing, so no Sobolev regularity of the input is required.

\subsection{History and significance}

Sogge introduced local smoothing in his study of circular maximal
operators and their variable-coefficient analogues \cite{Sogge1991}.
He proved a positive gain from time averaging in two spatial dimensions
and proposed the critical higher-dimensional formulation in
\cite[Section~4, especially (4.5)]{Sogge1991}. The conjecture measures
the improvement over the sharp fixed-time estimates of Peral and
Miyachi \cite{Peral1980,Miyachi1980}. In three dimensions it asks to
recover almost $1/p$ derivatives when $p\ge3$; at $p=3$, time
averaging should remove all but an arbitrarily small part of the
fixed-time loss of $1/3$. The connection with Fourier summation was
already part of this formulation: Sogge showed how critical smoothing
implies Bochner--Riesz bounds.

One approach separates the angular decomposition of the wave from the
geometry of its rays. Mockenhaupt--Seeger--Sogge developed square-function
estimates and their relation to local smoothing, including a framework
for Fourier integral operators satisfying cinematic curvature
\cite{MockenhauptSeegerSogge1992,MockenhauptSeegerSogge1993}.
A square function sums the squared magnitudes of the angular pieces at
each spacetime point before taking a norm. It therefore retains
information about where those pieces overlap. Geometric estimates for
thin tubes enter when this overlap is related to the initial data.
This approach gave partial smoothing gains and remained important
near the critical exponent.

Wolff obtained the conjectured gain for large $p$ in two spatial
dimensions by wave-packet localization and induction on scales
\cite{Wolff2000}. His decoupling approach estimates the sum of angular
pieces using their separate $L^p$ norms. {\L}aba--Wolff extended this
method to higher dimensions \cite{LabaWolff2002}, and
Garrig\'os--Seeger combined the induction scheme with bilinear
restriction estimates to improve cone plate-decomposition inequalities
\cite{GarrigosSeeger2009}. These results distinguished sharp derivative
gain from the full conjectured range of exponents. Bourgain--Demeter
then proved sharp cone decoupling in every dimension
\cite[Theorem~1.2]{BourgainDemeter2015}. Its local-smoothing consequence
gives the conjectured gain for $p\ge4$ in three spatial dimensions,
leaving the critical $p=3$ estimate beyond that range.

Beltran--Hickman--Sogge extended the decoupling method to variable
coefficients and proved sharp smoothing gains for wave equations on
compact Riemannian manifolds and a broader class of Fourier integral
operators \cite{BeltranHickmanSogge2020}. The geometric obstruction
depends on the setting. Minicozzi--Sogge constructed metrics with
concentrating families of geodesics \cite{MinicozziSogge1997}; in
spatial dimension three these examples require $p\ge10/3$ for the
full smoothing gain on general manifolds
\cite[Section~1.1]{BeltranHickmanSogge2020}. For the broader
cinematic-curvature class, the $p\ge4$ range proved by
Beltran--Hickman--Sogge is sharp in dimension three
\cite[Theorem~1.2 and Proposition~1.3]{BeltranHickmanSogge2020}.
These distinctions explain the Euclidean qualification in the
conjecture considered here.
In the broader setting of uniformly elliptic, time-independent wave
equations with real coefficients of bounded $C^{1,1}$ norm,
Rozendaal--Schippa obtain a diagonal $p=3$ estimate for $\alpha>2/9$
\cite[Corollary~1.2]{RozendaalSchippa2026}.

The two-dimensional completion illustrates the importance of retaining
spatial overlap information. There Bourgain--Demeter's theorem gives
the full smoothing gain for $p\ge6$; Guth--Wang--Zhang reached the
critical exponent $p=4$ through a sharp cone square-function estimate
\cite{GuthWangZhang2020}. Their stronger wave-envelope estimate tracks
how packets occupy regions at intermediate angular scales and exploits
sparse packet configurations. This has a direct role in the present
proof: we use the wave-envelope theorem as a lower-dimensional input,
while the three-dimensional argument must control the additional
direction and the changes of geometry between scales.

A further route uses transverse interactions between wave packets.
The multilinear estimates of Bennett--Carbery--Tao and the
multilinear-to-linear reduction of Bourgain--Guth
\cite{BennettCarberyTao2006,BourgainGuth2011} lead to a broad--narrow
strategy: separate contributions from many distinct directions from
those concentrated near a lower-dimensional direction space.
Guth's polynomial-partitioning method provides a weaker $k$-broad
estimate for the former, while the latter requires rescaling and
induction \cite{Guth2018}. Ou--Wang adapt this architecture to the
cone, combining broad estimates with narrow decoupling and Lorentz
rescaling \cite{OuWang2022}.

Gao--Liu--Miao--Xi use this broad--narrow architecture for local
smoothing. Rescaling changes the class of phases to be controlled;
their argument treats general positively curved cones and combines
broad estimates with narrow decoupling in scale-dependent phase
classes \cite{GaoLiuMiaoXi2023}. In three spatial dimensions they
obtain, at $p=3$, every Sobolev exponent $\alpha>1/9$, or equivalently
a smoothing gain below $2/9$
\cite[Corollary~1.3]{GaoLiuMiaoXi2023}.

Gan--He--Li--Wu combine refined decoupling with estimates for the local
$L^2$ energy and incidence geometry of wave packets
\cite{GanHeLiWu2026}. Their packet density measures energy in windows
compatible with Lorentz rescaling. A simultaneous broad and two-ends
reduction retains transverse contributions and controls concentration
in short portions of packets; in dimension three, slab coverings
sharpen the energy and incidence estimates together.
Theorem~1.3 of their January 2026 version gives the conjectured gain
for $p\ge10/3$, but does not include $p=3$. Interpolation of the
$p=10/3$ estimate with the $L^2$ energy bound gives $\alpha>1/12$
at $p=3$, improving the preceding $1/9$ threshold. The remaining
fixed positive loss is the gap addressed by Theorem~\ref{thm:main}.

Our proof also studies packet concentration across scales, but its
central quantity is the entropy of a finite phase-space record. The
records retain the changing shear frames as well as the time and
direction information. The proof uses two external quantitative inputs:
the planar Furstenberg theorem of Ren--Wang
\cite[Theorem~4.1]{RenWang2025} and the cone wave-envelope theorem of
Guth--Wang--Zhang \cite[Theorem~1.3]{GuthWangZhang2020}.
The former follows a line of discretized incidence and projection work
including Katz--Tao, Orponen--Shmerkin, and Shmerkin--Wang
\cite{KatzTao2001,OrponenShmerkin2023,ShmerkinWang2025}.
Section~\ref{sec:low} states these inputs and proves the specialized
probabilistic and Hilbert-valued consequences used below. The extremal
reductions and entropy inequalities then connect this lower-dimensional
information to the three-dimensional wave estimate.

\subsection{Structure of the argument}

We first seek the estimate on one frequency annulus
$|\xi|\asymp H^{-1}$: the physical half-wave should map $L^3$ into
spacetime $L^3$ with a loss of at most $H^{-\delta}$ for every
$\delta>0$. Summing these estimates gives Theorem~\ref{thm:main}.
On a fixed angular chart, null coordinates turn the light rays into
rays with velocities $(A,|A|^2)$, $A\in\R^2$. We analyze the wave
into packets and measure their energy by squared $L^2$ norms.

At the terminal time length $H$, the packet positions have width $H$
and their angular aperture has size one. Write $m_{j,Q}$ for the
packet energy assigned to a spatial $H$-cube $Q$ at the start of the
$j$th time interval. The synthesis estimate \eqref{eq:phys-synthesis}
controls the physical cubed $L^3$ norm by
$H^{-1/2}\sum_{j,Q}m_{j,Q}^{3/2}$, with an arbitrarily small power
loss and a negligible error. Thus the main task is to bound a cubic
sum of packet masses. Energy is copied by evolution into descendant
time intervals; those descendants are not orthogonal pieces of the
original input.

The intermediate scales require more flexible position and direction
windows. For the round model $(A,|A|^2)$, a window of angular radius
$r$ has position widths $\ell r$ and $\ell r^2$ at time length
$\ell$; a parameter $f\le1$ measures its relative thickness near a
hyperquadric. Thin round windows lead to the cylinder $(X,X^2,N)$.
There $r$ controls the curved base direction, while $f$ is the normal
velocity width, with normal position width $\ell f$ after a polynomial
shear. The models therefore use different test weights:
$r^4f^{1-\eta}$ in the round model and $r^{3+\eta}f^{1-\eta}$ in the
cylinder, where $0<\eta<1$. These windows and their weights are defined
in Definitions~\ref{cyl:def-test} and~\ref{rt:regularity}.

Suppose the initial packet measure has total energy $e$ and obeys
$\mu(Q)\le\kappa w(Q)$ for every test. At the final cut, let
$m_\lambda$ be the energy assigned to a test of weight $w_\lambda$.
The index $\lambda$ includes both its time interval and its test;
assignments may be fractional, but their sum at each interval is
bounded by that interval's row energy. The normalized cubic sum is
\begin{equation}\label{eq:intro-functional}
 \mathcal B_s=
 \frac{s^{-1/2}}{e\sqrt\kappa}
 \sum_\lambda\frac{m_\lambda^{3/2}}{\sqrt{w_\lambda}}.
\end{equation}
Here $s$ is the ratio between final and initial time lengths, and the
sum includes all final intervals. In the physical application $s=H$
and the terminal tests have weights comparable to one, recovering the
mass sum above. Theorem~\ref{thm:round} bounds $\mathcal B_s$ with a
scale loss that tends to zero with $\eta$ and the subsequent losses.
Section~\ref{sec:phys} verifies the required initial regularity on the
physical time slice and completes the frequency summation.

The model bounds are proved by contradiction. Suppose that the cubic
sum has a growth exponent larger than the permitted $\eta$-dependent
loss. In the cylinder model, an extremal reduction selects homogeneous
mass and weight data. Any ordinary regularity split that preserves the
extremal exponent must attain the corresponding energy and regularity
bounds at the level of exponents. Comparing the retained tests then
controls both their polynomial shears and their widths at successive
cuts. After a further reduction, the width exponents are affine
functions of depth. The lower-dimensional wave and incidence estimates
exclude the configurations in which both normalized shear costs vanish.

For the remaining configurations, we study a finite phase name at time
depth $a$ and horizontal resolution $v$. It records bins for time,
direction, base position, and the normal phase after subtracting a
recorded shear. From its normalized entropy we subtract the baseline
prescribed by regularity, the test weight, and the extremal mass data.
The resulting entropy excess is nonnegative. At specified boundary
resolutions, recovery of the name from a retained test label gives the
matching upper bound, so the excess vanishes there.
Conditional sampling supplies auxiliary probes, but the quantity
ultimately estimated remains the entropy of this one phase name.
Planar incidence and local projection estimates give directions in
which the excess must decrease when followed backwards from those
boundary points. Constructing a path that stays in the domain of the
applicable inequalities gives the contradiction to nonnegativity.
The admissible directions depend on the shear profile; unit-cost
boundaries and intervals with no angular entropy growth require
separate estimates.

To pass from the cylinder to the round velocity surface $(A,|A|^2)$,
$A\in\R^2$, we transfer thin round tests to cylindrical ones. If one
coordinate remains quantum, this transfer retains its one-dimensional
Schr\"odinger evolution. The remaining round configurations are excluded
by a separate entropy argument, first for classical rays and then for
waves. Finally, input regularity is verified on the original physical
time slice before summing the dyadic frequency estimates.

The paper is organized as follows.
Definitions~\ref{cyl:def-states}, \ref{cyl:def-test}, and
\ref{cyl:def-setup} specify the cylinder states, tests, and fixed
setups; Proposition~\ref{cyl:affine} gives the
extremal configurations used in the subsequent argument.
Section~\ref{sec:low} supplies the lower-dimensional estimates.
Definition~\ref{names:canonical-name} introduces the canonical names,
and Proposition~\ref{names:canonical-entropy} defines their entropy
excess. Section~\ref{int:section} derives the local inequalities,
which Section~\ref{act:section} uses in the proof of
Theorem~\ref{thm:cylinder}.
Definitions~\ref{rt:model} and~\ref{rt:regularity} specify the round
experiment and its tests. Proposition~\ref{rt:thin-exclusion} is the
thin-shape reduction used in Theorem~\ref{thm:round};
Section~\ref{sec:phys} then makes the passage to the physical wave.
Section~\ref{sec:bochner-riesz} derives the maximal and nonmaximal
Fourier-summation consequences and records their standard restriction
and Kakeya implications.

\begin{figure}[!htb]
\centering
\begin{tikzpicture}[
 node distance=7mm and 9mm,
 every node/.style={font=\small,align=center},
 result/.style={draw,inner sep=6pt,text width=43mm},
 >={Stealth[length=2mm]}]
 \node[result] (inputs) {Planar incidences\\Lower-dimensional waves};
 \node[result,right=of inputs] (cylinder) {Cylindrical packet estimates\\Theorem~\ref{thm:cylinder}};
 \node[result,below=of cylinder] (transfer) {Thin round tests\\Transfer to the cylinder};
 \node[result,left=of transfer] (round) {Round-cone estimates\\Theorem~\ref{thm:round}};
 \node[result,below=of round] (physical) {Physical input regularity\\Theorem~\ref{thm:main}};
 \draw[->] (inputs)--(cylinder);
 \draw[->] (cylinder)--(transfer);
 \draw[->] (transfer)--(round);
 \draw[->] (inputs)--(round);
 \draw[->] (round)--(physical);
\end{tikzpicture}
\caption{Dependencies of the main estimates. Arrows indicate which
estimates are used in subsequent arguments.}
\label{fig:dependencies}
\end{figure}

\begin{samepage}
The principal technical steps are the extremal reduction with compatible
shears, the phase-name calculus, and the uniform transfer of wave packets
on thin gaps. These intermediate results are formulated for the packet
models considered here, independently of the final frequency summation.
Throughout,
finite operator stacks and polynomial cutoff bounds are fixed before
taking a scale limit; Section~\ref{sec:framework} specifies the order of
all subsequent limits.
\par
\end{samepage}

\clearpage
\tableofcontents
\clearpage
